\documentclass[12pt]{article}
\usepackage[utf8]{inputenc}
\usepackage[T1]{fontenc}
\usepackage[margin=1in]{geometry}
\usepackage{setspace}
\usepackage{xcolor}
\usepackage[colorlinks=true, urlcolor=blue!60!black, linkcolor=blue!60!black]{hyperref}
\setlength{\parindent}{2em}
\setlength{\parskip}{1.25em}
\renewcommand{\baselinestretch}{1.0}
\title{\textit{Trump v. CASA, Inc.}, 606 U.S. 831 (2025)}
\author{Case Brief}
\date{}
\begin{document}
\maketitle

\section{Case Metadata}

\begin{itemize}
  \item \textbf{Case Name:} \textit{Trump v. CASA, Inc.} (consolidated with No.\ 24A885, \textit{Trump v. Washington}, and No.\ 24A886, \textit{Trump v. New Jersey})
  \item \textbf{Citation:} 606 U.S. 831 (2025); Docket No.\ 24A884 (on applications for partial stay)
  \item \textbf{Briefing Book Placement:} Chapter 1: Judicial Review and the Federal Courts (Equitable Remedies; Universal Injunctions; Judiciary Act of 1789)
  \item \textbf{Dates:} Applications Filed March 13, 2025; Argued May 15, 2025; Decided June 27, 2025
  \item \textbf{Vote:} 6--3 (granting applications for partial stays of three universal preliminary injunctions)
  \item \textbf{Opinion of the Court:} Justice Barrett, joined by Chief Justice Roberts and Justices Thomas, Alito, Gorsuch, and Kavanaugh
  \item \textbf{Concurring Opinions:} Justice Thomas (joined by Justice Gorsuch); Justice Alito (joined by Justice Thomas); Justice Kavanaugh
  \item \textbf{Dissenting Opinions:} Justice Sotomayor (joined by Justices Kagan and Jackson); Justice Jackson
  \item \textbf{Primary Sources:} \href{https://www.supremecourt.gov/opinions/24pdf/606us2r66_j426.pdf}{SupremeCourt.gov Official Page Proof}; \href{https://www.law.cornell.edu/supremecourt/text/24A884}{Cornell LII}; \href{https://supreme.justia.com/cases/federal/us/606/24a884/}{Justia U.S. Supreme Court Center}
\end{itemize}

\section{Facts and Procedural History}

On January 20, 2025, President Donald J. Trump issued Executive Order No.\ 14160, \textit{Protecting the Meaning and Value of American Citizenship}, 90 Fed.\ Reg.\ 8449 (2025), declaring that federal agencies would no longer issue or recognize U.S.\ citizenship documentation for children born in the United States whose mother was either unlawfully present or lawfully present only temporarily, and whose father was neither a U.S.\ citizen nor a lawful permanent resident at the time of birth (\href{https://www.law.cornell.edu/supremecourt/text/24A884}{Cornell LII}). The order included a 30-day ramp-up period for agencies to issue implementation guidance.

Individuals (including pregnant mothers), immigrant-rights organizations (including CASA, Inc.\ and Asylum Seeker Advocacy Project, representing hundreds of thousands of members nationwide), 22 States, the District of Columbia, and San Francisco filed three separate suits in federal district courts in Maryland (\textit{CASA, Inc.\ v. Trump}, 763 F. Supp.\ 3d 723 (D. Md.\ 2025)), the Western District of Washington (\textit{Washington v. Trump}, 765 F. Supp.\ 3d 1142 (W.D. Wash.\ 2025)), and Massachusetts (\textit{Doe v. Trump}, 766 F. Supp.\ 3d 266 (D. Mass.\ 2025)), alleging that Executive Order No.\ 14160 violated the Fourteenth Amendment's Citizenship Clause and \S\ 201 of the Nationality Act of 1940, 8 U.S.C.\ \S\ 1401. All three district courts concluded the order was likely unlawful and entered ``universal'' (nationwide) preliminary injunctions barring federal officials from enforcing it against anyone anywhere in the country.

The Fourth, Ninth, and First Circuits each denied the Government's stay motions (2025 WL 654902 (4th Cir.\ Feb.\ 28, 2025); 2025 WL 553485 (9th Cir.\ Feb.\ 19, 2025); 131 F.4th 27 (1st Cir.\ 2025)). On March 13, 2025, the Government filed three emergency applications in the Supreme Court seeking \textit{partial} stays---not defending the underlying constitutionality of Executive Order No.\ 14160, but arguing solely that the district courts lacked authority to enjoin enforcement beyond the parties. The Supreme Court consolidated the three applications and held oral argument on May 15, 2025 (\href{https://supreme.justia.com/cases/federal/us/606/24a884/}{Justia}).

\section{Issue(s) Presented}

Whether federal district courts possess equitable authority under the Judiciary Act of 1789 to issue ``universal injunctions'' that bar the Executive Branch from enforcing a challenged policy against nonparties beyond the relief necessary to afford complete relief to the plaintiffs with standing to sue. (The Court expressly did not decide whether Executive Order No.\ 14160 violates the Fourteenth Amendment or 8 U.S.C.\ \S\ 1401, nor whether Article III independently bars universal injunctions.)

\section{Holding}

No. In a 6--3 opinion authored by Justice Barrett, the Court held that universal injunctions likely exceed the equitable authority Congress conferred on federal courts in the Judiciary Act of 1789 (\href{https://www.law.cornell.edu/supremecourt/text/24A884}{Cornell LII}). The Court granted the Government's applications for partial stays, narrowing the three preliminary injunctions ``only to the extent that the injunctions are broader than necessary to provide complete relief to each plaintiff with standing to sue,'' and staying the bars on executive agencies developing public implementation guidance while ordering that \S\ 2 of the Executive Order not take effect for 30 days.

\section{Reasoning}

\subsection{The Judiciary Act of 1789 and Historical Equity}
Justice Barrett grounded the Court's ruling solely in statutory interpretation of \S\ 11 of the Judiciary Act of 1789 (1 Stat.\ 78), which grants federal courts jurisdiction over ``all suits \ldots\ in equity.'' Under \textit{Grupo Mexicano de Desarrollo, S.A.\ v. Alliance Bond Fund, Inc.}, 527 U.S.\ 308, 318--319 (1999), federal equitable authority is confined to remedies ``traditionally accorded by courts of equity''---specifically, the English High Court of Chancery at the founding. Surveying English and early American equity practice (\textit{Iveson v. Harris} (1802); \textit{Scott v. Donald}, 165 U.S.\ 107 (1897); \textit{Frothingham v. Mellon}, 262 U.S.\ 447 (1923); \textit{Perkins v. Lukens Steel Co.}, 310 U.S.\ 113 (1940); \textit{Doran v. Salem Inn, Inc.}, 422 U.S.\ 922 (1975)), the Court concluded that equitable decrees were traditionally party-specific and that universal injunctions were ``conspicuously nonexistent'' until the mid-twentieth century (e.g., \textit{Wirtz v. Baldor Elec.\ Co.}, 337 F.2d 518 (D.C. Cir.\ 1963)) (\href{https://www.law.cornell.edu/supremecourt/text/24A884}{Cornell LII}).

\subsection{Rejection of the Bill-of-Peace Analogy and Rule 23}
The Court rejected respondents' analogy to the English ``bill of peace.'' Unlike universal injunctions, bills of peace bound small, cohesive groups symmetrically (whether the group won or lost) and evolved into the modern class action governed by Federal Rule of Civil Procedure 23. Allowing single plaintiffs to obtain universal injunctions circumvents Rule 23's rigorous prerequisites (numerosity, commonality, typicality, adequacy) and creates asymmetric ``one-way'' litigation where plaintiffs need win only once while the Government must win everywhere. The Court explicitly reserved in footnote 10 whether the Administrative Procedure Act, 5 U.S.C.\ \S\ 706(2), authorizes universal vacatur of agency rules (\href{https://www.supremecourt.gov/opinions/24pdf/606us2r66_j426.pdf}{SupremeCourt.gov}).

\subsection{Complete Relief vs.\ Universal Relief and Stay Factors}
Distinguishing ``complete relief'' from ``universal relief,'' Justice Barrett explained that a court may issue an injunction that incidentally benefits nonparties when indivisible relief is necessary to redress the plaintiff's own injury (such as silencing a neighbor's nuisance or redrawing a gerrymandered map). Prohibiting enforcement against an individual pregnant plaintiff's child gives her complete relief without covering nonparties. For the state plaintiffs, however, the inquiry was ``more complicated'' because interstate migration and federally funded citizen-dependent benefits programs might render a patchwork injunction unworkable; the Court remanded for lower courts to determine whether narrower relief would suffice. Finally, under \textit{Nken v. Holder}, 556 U.S.\ 418 (2009), an overbroad injunction improperly intruding on a coordinate branch causes the Government irreparable harm even without reaching the underlying merits of Executive Order No.\ 14160 (\href{https://www.law.cornell.edu/supremecourt/text/24A884}{Cornell LII}).

\section{Concurring and Dissenting Opinions}

\begin{itemize}
  \item \textbf{Justice Thomas (Concurring, joined by Justice Gorsuch):} Emphasized that the complete-relief principle is a ``ceiling,'' not a ``mandate,'' and that indivisible nonparty-benefiting remedies are rare exceptions; courts must not smuggle universal injunctions back in under the guise of complete relief.
  \item \textbf{Justice Alito (Concurring, joined by Justice Thomas):} Warned of two potential loopholes that could gut \textit{CASA}: lax state third-party standing (allowing States to sue on behalf of all residents) and hasty certification of nationwide Rule 23 classes without ``scrupulous adherence'' to \textit{Wal-Mart Stores, Inc.\ v. Dukes}, 564 U.S.\ 338 (2011).
  \item \textbf{Justice Kavanaugh (Concurring):} Highlighted that \textit{CASA} addresses only ``the interim before the interim'' in district courts; plaintiffs may still seek Rule 23(b)(2) classwide relief or APA \S\ 706(2) set-aside orders, and the Supreme Court itself will remain the ultimate decider of nationally uniform interim legal status on emergency applications (\href{https://www.law.cornell.edu/supremecourt/text/24A884}{Cornell LII}).
  \item \textbf{Justice Sotomayor (Dissenting, joined by Justices Kagan and Jackson):} Condemned the Government's ``gamesmanship'' in seeking emergency relief without defending a patently unconstitutional order violating the Fourteenth Amendment and \textit{United States v. Wong Kim Ark}, 169 U.S.\ 649 (1898). She argued the Executive suffers no irreparable harm from being barred from breaking settled law, that bills of peace, taxpayer suits, and early twentieth-century cases (\textit{Pierce v. Society of Sisters}; \textit{West Virginia Bd.\ of Ed.\ v. Barnette}) establish historical pedigree for nonparty relief, and that even under the majority's complete-relief rule, the States' and organizations' nationwide injunctions were justified. She urged affected parents to file Rule 23(b)(2) class actions immediately.
  \item \textbf{Justice Jackson (Dissenting):} Argued that focusing on English Chancery analogies is ``legalese'' obscuring the rule-of-law premise that federal courts must be able to order the Executive to obey the law universally, warning that \textit{CASA} creates a ``zone of lawlessness'' disproportionately harming the poor and unrepresented. (Justice Barrett's majority sharply countered that Jackson's view embraced an ``imperial Judiciary.'')
\end{itemize}

\section{Key Quotations}

\begin{quote}
``But federal courts do not exercise general oversight of the Executive Branch; they resolve cases and controversies consistent with the authority Congress has given them. When a court concludes that the Executive Branch has acted unlawfully, the answer is not for the court to exceed its power, too.'' --- Justice Barrett, Opinion of the Court (\href{https://www.law.cornell.edu/supremecourt/text/24A884}{Cornell LII}).
\end{quote}

\begin{quote}
````Complete relief'' is not synonymous with ``universal relief.'' It is a narrower concept: The equitable tradition has long embraced the rule that courts generally ``may administer complete relief between the parties.'''' --- Justice Barrett, Opinion of the Court (\href{https://www.law.cornell.edu/supremecourt/text/24A884}{Cornell LII}).
\end{quote}

\begin{quote}
``No right is safe in the new legal regime the Court creates. ... That holding renders constitutional guarantees meaningful in name only for any individuals who are not parties to a lawsuit.'' --- Justice Sotomayor, dissenting (\href{https://www.law.cornell.edu/supremecourt/text/24A884}{Cornell LII}).
\end{quote}

\begin{quote}
``The Court's decision to permit the Executive to violate the Constitution with respect to anyone who has not yet sued is an existential threat to the rule of law.'' --- Justice Jackson, dissenting (\href{https://www.law.cornell.edu/supremecourt/text/24A884}{Cornell LII}).
\end{quote}

\section{Doctrinal Significance and Modern Legacy (Through 2026)}

\textit{Trump v. CASA, Inc.}\ reshaped federal remedial practice while leaving open three major pathways to broad relief---producing a nuanced first-year legacy documented in the Congressional Research Service's July 7, 2026 study (\href{https://www.everycrsreport.com/reports/R49022.html}{CRS Report R49022}):

\begin{enumerate}
  \item \textbf{Pivot to Nationwide Rule 23(b)(2) Class Actions and the Merits Demise of EO 14160:} Following Justice Sotomayor's and Justice Kavanaugh's roadmaps, challengers swiftly filed class actions, securing a classwide preliminary injunction on remand in \textit{CASA, Inc.\ v. Trump}, 793 F. Supp.\ 3d 687 (D. Md.\ 2025), and provisional nationwide class certification in \textit{Barbara v. Trump}, 2025 WL 1904338 (D.N.H.\ July 10, 2025) (\href{https://houstonlawreview.org/article/171991-from-universal-injunctions-to-nationwide-classes-what-_barbara-v-trump_-reveals-about-_casa_-s-remedial-framework.pdf}{Houston Law Review}). Granting certiorari before judgment in \textit{Trump v. Barbara}, No.\ 25-365, the Supreme Court resolved the underlying merits on June 30, 2026, holding 6--3 that Executive Order No.\ 14160 violates the Fourteenth Amendment's Citizenship Clause (\href{https://www.supremecourt.gov/opinions/25pdf/25-365_4hdj.pdf}{SupremeCourt.gov}). Appellate courts have nonetheless policed Rule 23 boundaries, with the Eighth Circuit staying an uncertified-class injunction as ``a universal injunction by another name'' in \textit{Tincher v. Noem}, 164 F.4th 1097 (8th Cir.\ 2026).
  \item \textbf{State-Plaintiff Complete-Relief Survival:} On remand from \textit{CASA}, both the Ninth Circuit (\textit{Washington v. Trump}, 145 F.4th 1013 (9th Cir.\ 2025)) and First Circuit (\textit{Doe v. Trump}, 157 F.4th 36 (1st Cir.\ 2025)) affirmed nationwide injunctions as necessary to give plaintiff States complete relief. Overall, CRS tracked a slowdown from 34 nationwide injunctions in the five months pre-\textit{CASA} to 16 in the full year post-\textit{CASA} (\href{https://www.everycrsreport.com/reports/R49022.html}{CRS Report R49022}).
  \item \textbf{APA \S\ 706(2) Vacatur Carveout:} Lower courts through 2026 have uniformly refused to extend \textit{CASA}'s Judiciary Act holding to universal stays or vacatur of agency rules under 5 U.S.C.\ \S\ 706(2) (e.g., \textit{Louisiana v. FDA}, 175 F.4th 310, 323 (5th Cir.\ 2026); \textit{City of Columbus v. Kennedy}, 796 F. Supp.\ 3d 123 (D. Md.\ 2025)), leaving statutory APA vacatur as the next major remedial battleground.
\end{enumerate}

\section{Pedagogical Critiques and Exam Hypotheticals}

\subsection{Pedagogical Critiques}
\begin{itemize}
  \item \textbf{Statutory Originalism vs.\ Dynamic Equity:} Critics question whether freezing federal equity to 1789 English Chancery practice under \textit{Grupo Mexicano} makes sense when modern administrative and presidential governance---and officer suits under \textit{Ex parte Young}---had no founding-era Chancery twin.
  \item \textbf{Formalism vs.\ Practical Workarounds:} Because \textit{CASA} rests on the Judiciary Act of 1789 rather than Article III, Congress may alter the rule by statute, and plaintiffs have largely replicated nationwide relief via Rule 23(b)(2) classes, State-plaintiff complete-relief showings, and APA \S\ 706(2) vacatur---raising the question whether \textit{CASA} merely shifted procedural friction onto the Supreme Court's emergency docket.
\end{itemize}

\subsection{Exam Hypotheticals}
\begin{enumerate}
  \item \textbf{Hypothetical 1 (Individual Challenge to Presidential Executive Order):} The President issues an executive order directing the Treasury to impose a 10\% surcharge on all international wire transfers sent by U.S.\ residents. Two individual plaintiffs sue in federal district court, alleging the order violates Article I's Taxing Clause, without moving for Rule 23 class certification or naming any agency rule under the APA. Finding the order patently unconstitutional, the district court enters a universal preliminary injunction barring enforcement nationwide. \textit{Analysis:} Under \textit{CASA}, the universal scope must be stayed. Enjoining the Treasury from collecting the surcharge on the two named plaintiffs' transfers affords them complete relief; extending protection to nonparties exceeds the district court's equitable authority under the Judiciary Act of 1789.
  \item \textbf{Hypothetical 2 (``Opposite Side'' Contrast --- APA Vacatur and Certified Rule 23(b)(2) Class Action):} Suppose instead that (a) a trade association sues under the Administrative Procedure Act, 5 U.S.C.\ \S\ 706(2), to ``hold unlawful and set aside'' a final Treasury regulation implementing the surcharge, OR (b) the individual plaintiffs satisfy Rule 23(a) and 23(b)(2) and obtain certification of a nationwide class of all U.S.\ senders subject to the surcharge before seeking a classwide injunction. \textit{Analysis:} Unlike \textit{CASA}, neither remedy is barred by \textit{CASA}'s holding. Footnote 10 explicitly carved out APA \S\ 706(2) vacatur (which lower courts through 2026 continue to apply universally), and the majority itself identified a properly certified Rule 23(b)(2) class action as the lawful modern descendant of the English bill of peace---provided the district court conducts the ``rigorous analysis'' demanded by \textit{Wal-Mart} and Justice Alito's concurrence.
\end{enumerate}

\end{document}
